\cgeentry{cge_iis}{An irreducible infeasible set of bounds and constraints}

\cgesec{Description}

\cmd{cge_iis} (from version 0.36.0) looks for the reason a scenario cannot be
solved among its bounds (\manref{cge_bound}) and constraints
(\manref{cge_constraint}, and the inequalities of \manref{cge_equation}). It
finds an irreducible infeasible set, as the IIS of CPLEX and Gurobi does for
a linear programme: bounds and constraints that together leave the scenario
without a solution, while it solves without any one of them. A scenario
that fails with dozens of bounds and constraints is thereby reduced to the
few that conflict.

\cgesec{Quick start}

The set among the bounds and constraints of the current scenario
\begin{stcode}
cge_iis
\end{stcode}
The same, with the outcome of each step
\begin{stcode}
cge_iis, trace
\end{stcode}

\cgesec{Syntax}

\begin{stsyntax}
\cmdb{cge_iis} [\texttt{,} \optarg{tol}{erance}{\#} \optarg{iter}{ate}{\#} \opt{}{trace}]
\end{stsyntax}

\begin{opttable}
\optgroup{Main}
\optarg{tol}{erance}{\#} & tolerance of each solve; default that of \manref{cge_options}\\
\optarg{iter}{ate}{\#} & iteration limit of each solve; default that of \manref{cge_options}\\
\opt{}{trace} & show the outcome of each step\\
\end{opttable}

\cgesec{Remarks and examples}

The search is the deletion filter of \citet{chinneck1991}. The scenario must
fail to solve with every bound and constraint, and solve without them all;
each in turn is then left out for good if the scenario still fails without
it. What remains is the set. The search takes one solve for each bound and
constraint, and two more. A failure to solve stands for infeasibility, so the
set is that of the solver as much as of the model \citep{chinneck2008}: a
scenario that the solver fails on for another reason gives a set as well.

Each solve is of the square system of the scenario, as \manref{cge_solve}
without \opt{}{minimize()} or \opt{}{maximize()} makes it, and at the first of
three levels at which the scenario without its bounds and constraints
solves: Newton's method from the base; with the fallbacks of
\manref{cge_solve} (Newton's method with central differences, then Powell's
hybrid method); and with them in 4 and then 16 steps. Most scenarios solve
at the first level, where a solve that fails stops at once: the search then
takes seconds where the fallbacks would take minutes.

The equations themselves are not candidates: without one, the system would
not be square. When the scenario cannot be solved even without its bounds
and constraints, \cmd{cge_iis} stops with an error; \manref{cge_listing}
shows where the solve stops, and \manref{cge_structure} whether the system
is structurally singular.

The SAM is mapped and calibrated:
\stlog{iis_setup}

\cgesubsec{Example 1: a bound among bounds that do not bind}

A tariff removal raises imports. With imports of manufactures held at
most at their base, and bounds on other imports and on the composite
prices that the solution does not reach, the scenario has no solution. The
set is the one bound that cannot hold:
\stlog{iis_bounds}

\cgesubsec{Example 2: a constraint and a bound together}

A quota holds imports of manufactures below 98\% of their base, with the
rent of the quota adjusting; the scenario solves, with a rent of 0.0974.
Capped at half that, the rent cannot rise far enough. Neither the
inequality nor the cap is infeasible alone: the set is the two together.
\stlog{iis_quota}

\cgesec{Stored results}

\cmd{cge_iis} stores the following in \cmd{r()}:
\begin{storedres}
\resgroup{Scalars}
\res{r(N_candidates)} & bounds and constraints of the scenario\\
\res{r(N_iis)} & bounds and constraints in the set (0 if the scenario solves)\\
\res{r(N_solves)} & solves of the search\\
\resgroup{Macros}
\res{r(iis)} & the bounds and constraints of the set\\
\end{storedres}

\cgesec{Methods and formulas}

With candidates $C = \{c_1, \ldots, c_n\}$, the bounds in the order of
\cmd{cge_bound, list} and then the constraints, and $F(S)$ true when the
scenario with the candidates in $S$ cannot be solved, the search checks
$F(\emptyset)$ false and $F(C)$ true, and then, with $S = C$, for $k = 1,
\ldots, n$, sets $S = S \setminus \{c_k\}$ where $F(S \setminus \{c_k\})$
holds. The final $S$ has $F(S)$ true, and $F(S \setminus \{c\})$ false for
each $c$ in it: when $c$ was examined, the set without it solved, and every
set since is a subset of that one with $c$ \citep{chinneck1991}. A bound on
a variable that the closure fixes outside it fails without a solve.

\cgesec{Also see}

\begin{alsosee}
\seeentry{cge_bound}{Attributes of variables: bounds, scales, starting levels and fixed values}
\seeentry{cge_constraint}{Constraints on variables, with an adjusting variable}
\seeentry{cge_listing}{Equation listing, infeasibility report and solution listing}
\seeentry{cge_structure}{Structure of the system: structural singularity and blocks}
\seeentry{cge_solve}{Solve the model}
\end{alsosee}
